\documentclass[12pt,a4paper]{article}

%=========================================================
%                   ENCODAGE ET POLICES
%=========================================================
\usepackage[utf8]{inputenc}
\usepackage[T1]{fontenc}
\usepackage[french]{babel}
\usepackage{lmodern}

%=========================================================
%                     MATHÉMATIQUES
%=========================================================
\usepackage{amsmath,amssymb,amsfonts,mathrsfs}
\usepackage{tkz-tab}
\everymath{\displaystyle}

%=========================================================
%                     MISE EN PAGE
%=========================================================
\usepackage[left=1.5cm,right=1.5cm,top=1.5cm,bottom=1.8cm]{geometry}
\usepackage{multicol}
\usepackage{float}
\usepackage{array,multirow}
\usepackage{enumitem}
\usepackage{color, soul}

%=========================================================
%                GRAPHISMES ET COULEURS
%=========================================================
\usepackage{tikz}
\usetikzlibrary{arrows.meta,calc}
\usepackage{pgfplots}
\pgfplotsset{compat=1.15}

\usepackage[most]{tcolorbox}
\tcbuselibrary{skins,hooks}

%=========================================================
%                  LIENS ET ARRIÈRE-PLAN
%=========================================================
\usepackage[colorlinks=true,linkcolor=blue,urlcolor=blue,citecolor=blue]{hyperref}
\usepackage{eso-pic}

%=========================================================
%                   COULEURS ET STYLES
%=========================================================
\definecolor{myred}{RGB}{180, 40, 40}
\definecolor{myblue}{RGB}{20, 50, 120}
\definecolor{mygreen}{RGB}{30, 120, 40}

%=========================================================
%                       FILIGRANE
%=========================================================
\AddToShipoutPicture{
    \AtTextCenter{
        \makebox[0pt]{
            \rotatebox{45}{
                \textcolor[gray]{0.92}{
                    \fontsize{5cm}{5cm}\selectfont PGB
                }
            }
        }
    }
}

\begin{document}

% En-tête
\begin{center}
    \Large\textbf{\underline{\textcolor{myred}{LE THÉORÈME DE THALÈS}}}\\[0.2cm]
    \normalsize\textbf{Prof : M. BA} \hfill \textbf{Classe : 3 ème}\\[0.1cm]
    \textbf{Durée : 6 heures} \hfill \textbf{Année scolaire : 2025 -- 2026}
\end{center}
\hrule
\vspace{0.3cm}

%=========================================================
%    I. La racine carrée d'un nombre réel positif
%=========================================================
\section*{\color{myred}I. Le théorème de Thalès direct}
\subsubsection*{\color{myred} Activité 1}

\subsection*{\color{myred}1. Énoncé du théorème}

\subsection*{\color{myred}Remarque}

\subsection*{\color{myred} 2. Calcul d'une longueur}
\subsubsection*{\color{myred} Exemple 1}
\subsection*{\color{myred} Application 1}
\subsection*{\color{myred} Correction 1}

\subsection*{\color{myred} 3. Montrer que deux droites ne sont pas parallèles}
\subsubsection*{\color{myred} Exemple 2}
\subsection*{\color{myred} Application 2}
\subsection*{\color{myred} Correction 2}

\section*{\color{myred}II. La réciproque du théorème de Thalès}
\subsection*{\color{myred} Activité 2}
\subsection*{\color{myred} 1. Énoncé du théorème}

\subsection*{\color{myred} 2. Montrer que deux droites sont parallèles}
\subsection*{\color{myred} Activité 3}
\subsubsection*{\color{myred} Attention}
\subsubsection*{\color{myred} Exemple 3}
\subsection*{\color{myred} Application 3}
\subsection*{\color{myred} Correction 3}

\end{document}